% Wstawki do roznych sekcji skryptu "Geometria i topologia rozniczkowa".
% Korzystaja z istniejacych srodowisk przyklad, makr \dd, \R
% i pakietow tikz oraz pgfplots (obecnych w preambule skryptu).
% Konwencja: R(X,Y)Z=nabla_X nabla_Y Z - nabla_Y nabla_X Z - nabla_[X,Y]Z.

% ================================================================
% CODEX / INSTRUKCJA INTEGRACJI Z GLOWNYM SKRYPTEM
% ================================================================
% Ten plik zawiera 15 przykladow i definicje izomorfizmow muzycznych
% do wstawienia w ISTNIEJACE sekcje skryptu "Geometria i topologia rozniczkowa".
% NIE wstawiaj calego pliku poleceniem \\input na koncu skryptu:
% porzadek definicji i odsylaczy ma pozostac logiczny.
%
% 0) def:izomorfizmy-muzyczne + prz:izomorfizmy-muzyczne:
%    rozdzial "Tensor metryczny, dlugosc i miara", sekcja
%    \label{sec:metryka-definicja}; po wzorze \label{eq:metryka-macierz}
%    i akapicie wyjasniajacym zapis g_ij dx^i dx^j, PRZED
%    \begin{stwierdzenie}[Zmiana wspolrzednych].
%    W skrypcie z 5 pazdziernika 2026 nie ma definicji
%    \flat ani \sharp: zachowaj definicje i przyklad.
% 1) prz:poincare-czas-urojony -- rozdzial o metryce:
%    znajdz \\label{def:metryka-lorentzowska}; wstaw po akapicie
%    zaczynajacym sie "Tu slowo metryka oznacza tensor", a przed
%    "W ponizszych dwoch przykladach uzywamy jednostek geometrycznych".
% 2) prz:flrw-metryka -- ten sam rozdzial: po akapicie
%    "W obu przykladach zapisalismy konkretne tensory...", a przed
%    \\section{Dlugosc, energia i odleglosc} (naglowek z polskimi znakami).
% 3) prz:flrw-christoffel -- sekcja koneksji Leviego-Civity:
%    bezposrednio po wzorze \\label{eq:christoffel-metryka} i jego objasnieniu.
% 4) prz:poincare-geodezyjne: sekcja geodezyjnych, PO
%    wyprowadzeniu rownania geodezyjnej \label{eq:geodezyjna-ode};
%    na przyklad przed sekcja o Schwarzschildzie i Kerrze.
%    prz:flrw-geodezyjne i prz:flrw-jacobi -- rozdzial geodezyjnych:
%    po \\section{Schwarzschild i Kerr: geodezyjne czasoprzestrzeni}
%    (lub pod koniec tej sekcji), w podanej kolejnosci.
% 5) prz:poincare-krzywizna: po definicji Riemanna, Ricciego,
%    krzywizny skalarnej i wspolrzednych tensora Riemanna,
%    \label{eq:riemann-wspolrzedne}; przed prz:flrw-krzywizna.
%    prz:flrw-krzywizna -- rozdzial o krzywiznie, po
%    \\label{eq:riemann-wspolrzedne} i definicji Ricciego i skalara,
%    przed \\section{Schwarzschild i Kerr: proznia nie oznacza plaskosci}.
% 6) prz:flrw-friedmann, prz:flrw-rozwiazanie-ode,
%    prz:flrw-konforemny, prz:flrw-osobliwosc,
%    prz:flrw-czas-zespolony, prz:flrw-milne -- pod koniec
%    rozdzialu o krzywiznie, po omowieniu Schwarzschilda i Kerra.
%    Zachowaj DOKLADNIE ta kolejnosc, bo uzyto wspolnych etykiet.
%
% WAZNE: nie duplikuj etykiet eq:*, prz:* ani fig:*;
% nie tworz nowych srodowisk i nie dopisuj pakietow, jesli sa.
% Preambula glownego skryptu zawiera tikz, pgfplots, hyperref,
% kolory manifoldblue, manifoldred, manifoldgold i makra \\dd, \\R.
% W tej bibliotece c=1 w FLRW; c pozostaje jawne u Poincarego.
% W miejscach wstawienia usun tylko komentarze CODEX, nie wzory.
% Nie dodawaj nowego rozdzialu z matematyka Poincarego;
% rozdziel przyklady miedzy METRYKE, GEODEZYJNE i KRZYWIZNE.
% Nie wstawiaj notacji flat/sharp ponownie do rozdzialu tensorow,
% jezeli jest juz w wersji roboczej; sprawdz faktyczny plik.
% Obliczenia przy t=0 dotycza MINKOWSKIEGO, nie FLRW!
% ================================================================

% ===== CODEX: DEFINICJA + PRZYKLAD; ZARAZ PO eq:metryka-macierz =====
% ===== Rozdzial "Tensor metryczny, dlugosc i miara" =====
% Nie jest to deklaracja makr w preambule: \flat i \sharp
% sa standardowymi symbolami matematycznymi w LaTeX.

\begin{definicja}[Izomorfizmy muzyczne: \texorpdfstring{$\flat$}{flat} i
\texorpdfstring{$\sharp$}{sharp}]
\label{def:izomorfizmy-muzyczne}
Niech $g$ będzie metryką riemannowską na $M$.
Jej niezdegenerowanie pozwala w każdym punkcie utożsamić
wektory styczne z kowektorami, definiując
\begin{equation}\label{eq:muzyczne-flat-sharp}
 \flat_g\colon TM\longrightarrow T^*M,\qquad
 v\longmapsto v^\flat:=g(v,\cdot),\qquad
 \sharp_g:=\flat_g^{-1}\colon T^*M\longrightarrow TM,
 \quad \alpha\longmapsto\alpha^\sharp.
\end{equation}
Wektor $\alpha^\sharp$ jest więc jednoznacznie określony
warunkiem $g(\alpha^\sharp,w)=\alpha(w)$ dla każdego $w$.
Nazwy \emph{obniżanie} i \emph{podnoszenie indeksu}
odnoszą się do współrzędnych tych odwzorowań.
Jeżeli $g=g_{ij}\dd x^i\dd x^j$ oraz
$(g^{ij})=(g_{ij})^{-1}$, to dla
$X=X^i\partial_i$ i $\alpha=\alpha_i\dd x^i$ mamy
\begin{equation}\label{eq:muzyczne-wspolrzedne}
 X^\flat=g_{ij}X^j\dd x^i,\qquad
 \alpha^\sharp=g^{ij}\alpha_j\partial_i,
 \qquad (X^\flat)^\sharp=X,\quad
 (\alpha^\sharp)^\flat=\alpha.
\end{equation}
Tak samo podnosimy i obniżamy wybrane indeksy tensorów.
Dla tensora $S=S_{ij}\dd x^i\otimes\dd x^j$ otrzymujemy
$(S^{\sharp_1})^i{}_j=g^{ik}S_{kj}$ i
$(S^{\sharp_1\sharp_2})^{ij}=g^{ik}g^{j\ell}S_{k\ell}$;
należy wskazać, który indeks podnosimy.
W szczególności $(g^{\sharp_1})^i{}_j=\delta^i_j$.
Definicja pozostaje poprawna dla każdej niezdegenerowanej
metryki pseudoriemannowskiej, także lorentzowskiej;
dodatnia określoność nie jest tu potrzebna.
\end{definicja}

\begin{przyklad}[obniżanie i podnoszenie indeksów]
\label{prz:izomorfizmy-muzyczne}
Na $\R^2$ rozważmy $g=\dd x^2+4\dd y^2$.
Dla $X=\partial_x+3\partial_y$ oraz
$\alpha=2\dd x+8\dd y$ zachodzi
$X^\flat=\dd x+12\dd y$ i
$\alpha^\sharp=2\partial_x+2\partial_y$.
Istotnie, macierz metryki to $G=\mathrm{diag}(1,4)$,
a jej odwrotność wynosi $G^{-1}=\mathrm{diag}(1,1/4)$.
Widać, że $X^\flat$ nie jest jedynie zmianą sposobu
zapisania tych samych współczynników: zależy od metryki.

Dla metryki Minkowskiego
$g_M=-c^2\dd t^2+\dd x^2+\dd y^2+\dd z^2$ mamy natomiast
$(\partial_t)^\flat=-c^2\dd t$ oraz
$(\dd t)^\sharp=-c^{-2}\partial_t$.
Znak minus wynika z sygnatury lorentzowskiej.
W szczególności, jeżeli $U=\partial_t$ i $c=1$,
to $U^\flat=-\dd t$; tego zapisu użyjemy
później w tensorze energii i pędu płynu doskonałego.
\end{przyklad}

% ===== CODEX: PRZYKLAD HISTORYCZNY; PO def:metryka-lorentzowska =====

\begin{przyklad}[urojona współrzędna czasowa Poincarégo]
\label{prz:poincare-czas-urojony}
W pracy \emph{Sur la dynamique de l'électron}, przedłożonej
w 1905 roku i opublikowanej w 1906 roku, Henri Poincaré
interpretował transformacje Lorentza jako obroty
w czterowymiarowej przestrzeni z urojoną współrzędną
czasową. Nie jest to jeszcze ogólna teoria względności:
rozważana geometria jest płaska.

Na $M=\R^4$, we współrzędnych $(t,x,y,z)$, rozważmy
metrykę Minkowskiego
\begin{equation}\label{eq:poincare-minkowski}
 g_M=-c^2\dd t^2+\dd x^2+\dd y^2+\dd z^2.
\end{equation}
Ma ona sygnaturę $(-,+,+,+)$. W dalszych przykładach
\ref{prz:poincare-geodezyjne} i~\ref{prz:poincare-krzywizna}
obliczymy jej koneksję, geodezyjne i pełny tensor
krzywizny, a następnie sprawdzimy regularność
hiperpowierzchni $t=0$.

Wprowadźmy $\zeta=\mathrm{i}ct$, gdzie $\mathrm{i}^2=-1$.
Wtedy $\dd\zeta=\mathrm{i}c\,\dd t$ i formalnie
\begin{equation}\label{eq:poincare-complex-metric}
 g_M=\dd x^2+\dd y^2+\dd z^2+\dd\zeta^2.
\end{equation}
Nie jest to dodatnio określona metryka na rzeczywistej
czasoprzestrzeni: $\zeta\in\mathrm{i}\R$, nie zaś
$\R$. Ściślej, na $\C^4$ rozważmy \emph{zespoloną
symetryczną formę dwuliniową} (nie hermitowską)
$h=\sum_{j=1}^4\dd z_j^2$ i odwzorowanie
$F\colon\R^4\to\C^4$ dane przez
$F(t,x,y,z)=(x,y,z,\mathrm{i}ct)$.
Wówczas $F^*h=g_M$. Jest to więc algebraiczne
przedstawienie tej samej rzeczywistej metryki
lorentzowskiej za pomocą formy zespolonej.

Dla prędkości $|v|<c$ połóżmy $\beta=v/c$ oraz
$\gamma=(1-\beta^2)^{-1/2}$. Transformacja Lorentza
w kierunku osi $x$ jest dana wzorami
$x'=\gamma(x-vt)$,
$t'=\gamma(t-vx/c^2)$, $y'=y$, $z'=z$.
W zmiennych $(x,\zeta)$ przyjmuje postać
\begin{equation}\label{eq:poincare-complex-boost}
 \begin{pmatrix}x'\\\zeta'\end{pmatrix}
 =\gamma\begin{pmatrix}
  1&\mathrm{i}\beta\\-\mathrm{i}\beta&1
 \end{pmatrix}
 \begin{pmatrix}x\\\zeta\end{pmatrix}.
\end{equation}
Macierz $A$ tego przekształcenia spełnia
$A^TA=I$ oraz $\det A=1$, czyli $A\in SO(2,\C)$.
Jeśli $\beta=\tanh\varphi$, to
$\gamma=\cosh\varphi$ i $\gamma\beta=\sinh\varphi$:
uzyskujemy obrotową macierz o \emph{zespolonym}
kącie $-\mathrm{i}\varphi$. W rzeczywistych
współrzędnych jest to obrót hiperboliczny, zachowujący
$x^2+y^2+z^2-c^2t^2$ oraz stożek świetlny,
a nie zwykły obrót euklidesowy.

\begin{figure}[htbp]
\centering
\begin{tikzpicture}[>=Stealth,x=1.04cm,y=1.04cm]
 \draw[->] (-2.65,0)--(2.70,0) node[right] {$x$};
 \draw[->] (0,-.25)--(0,3.22) node[above] {$ct$};
 % Stożek świetlny w przekroju przestrzenno-czasowym
 \draw[gray!70,densely dashed] (-2.5,2.5)--(0,0)--(2.5,2.5);
 \node[gray!75!black,above left] at (-1.45,1.45) {$ct=-x$};
 \node[gray!75!black,above right] at (1.35,1.35) {$ct=x$};
 % Hiperbola o stałym interwale
 \draw[manifoldblue,very thick,domain=-2.4:2.4,samples=75]
   plot (\x,{sqrt(1+\x*\x)});
 \node[manifoldblue,fill=white,inner sep=1.5pt] at (-1.8,2.9)
   {$c^2t^2-x^2=1$};
 % Osie układu primowanego dla beta=0.4
 \draw[manifoldgold!85!black,thick,->] (0,0)--(.93,2.33)
   node[above right] {$ct'$};
 \draw[manifoldred!85!black,thick,->] (0,0)--(2.36,.94)
   node[right] {$x'$};
 \fill (0,0) circle (1.4pt);
\end{tikzpicture}
\caption{Przekrój czasoprzestrzeni Minkowskiego $(x,ct)$.
Transformacja Lorentza zachowuje proste zerowe
$ct=\pm x$ i hiperbole stałego interwału. Osie
$(x',ct')$ układu poruszającego się ($\beta=0{,}4$)
nie powstają przez obrót euklidesowy.}
\label{fig:poincare-boost}
\end{figure}

Podstawienie $\zeta=\mathrm{i}ct$ nie jest zatem
rzeczywistą zmianą współrzędnych ani procedurą
usuwania osobliwości. Prawo bezwładności Sylvestera
wyklucza zamianę sygnatury lorentzowskiej na
riemannowską przez rzeczywistą zmianę bazy.
\end{przyklad}


% ===== CODEX: PRZYKLAD RACHUNKOWY; ROZDZIAL GEODEZYJNYCH =====
% Wstaw po zdefiniowaniu symboli Christoffela i rownania
% geodezyjnej eq:geodezyjna-ode, a przed przykladami FLRW.

\begin{przyklad}[geodezyjne i symbole Christoffela u Poincarégo]
\label{prz:poincare-geodezyjne}
Wróćmy do metryki Minkowskiego~\eqref{eq:poincare-minkowski}.
Przyjmijmy $x^0=t$ oraz $(x^1,x^2,x^3)=(x,y,z)$.
Macierz metryki i macierz odwrotna są stałe:
\[
 (g_{\mu\nu})=\operatorname{diag}(-c^2,1,1,1),\qquad
 (g^{\mu\nu})=\operatorname{diag}(-c^{-2},1,1,1).
\]
Dla dowolnych $\mu,\nu,\rho$ mamy
$\partial_\rho g_{\mu\nu}=0$, a więc wzór
na koneksję Leviego-Civity daje
\begin{equation}\label{eq:poincare-christoffel}
 \Gamma^\lambda_{\mu\nu}
 =\frac12g^{\lambda\rho}
 (\partial_\mu g_{\rho\nu}+\partial_\nu g_{\rho\mu}
 -\partial_\rho g_{\mu\nu})=0.
\end{equation}
Równanie geodezyjnej przyjmuje zatem postać
$(x^\mu)''=0$. Jest to układ czterech niezależnych
liniowych równań różniczkowych zwyczajnych drugiego
rzędu, którego wszystkie rozwiązania mają postać
\begin{equation}\label{eq:poincare-geodezyjne}
 \gamma(s)=(t_0+u^0s,\,x_0+u^1s,\,
 y_0+u^2s,\,z_0+u^3s),\qquad s\in\R.
\end{equation}
Wektor styczny $u=\gamma'$ jest stały i
$g_M(u,u)=-c^2(u^0)^2+(u^1)^2+(u^2)^2+(u^3)^2$.
Krzywe są odpowiednio czasowe, zerowe lub przestrzenne,
gdy to wyrażenie jest ujemne, zerowe lub dodatnie.
Dla $u^0\neq0$ geodezyjna przecina hiperpowierzchnię
$t=0$ przy $s_0=-t_0/u^0$ i przedłuża się na obie
strony tego punktu bez zmiany wzoru.
Każda geodezyjna ma dziedzinę całe $\R$,
więc $(\R^4,g_M)$ jest geodezyjnie zupełna.

Podstawienie Poincarégo $\zeta=\mathrm{i}ct$ przenosi
te proste na $F\circ\gamma$, gdzie
$\zeta(s)=\mathrm{i}c(t_0+u^0s)$ i $\zeta''(s)=0$.
Ograniczenie $\zeta\in\mathrm{i}\R$ jest istotne:
nie pojawia się dodatkowy rzeczywisty kierunek czasowy.

\begin{figure}[htbp]
\centering
\begin{tikzpicture}[>=Stealth,x=1.05cm,y=.87cm]
 \draw[->] (-2.7,0)--(2.8,0) node[right] {$x$};
 \draw[->] (0,-2.45)--(0,2.55) node[above] {$ct$};
 \draw[gray!60,densely dashed] (-2.4,-2.4)--(2.4,2.4);
 \draw[gray!60,densely dashed] (-2.4,2.4)--(2.4,-2.4);
 \draw[manifoldblue,very thick,->] (-.9,-2.2)--(.9,2.2);
 \draw[manifoldgold!90!black,very thick,->] (1.4,-2.12)--(.4,2.25);
 \fill[manifoldblue] (0,0) circle (2.2pt);
 \fill[manifoldgold!90!black] (.913,0) circle (2.2pt);
 \node[manifoldblue,above left] at (1.5,1.75) {$\gamma_1$};
 \node[manifoldgold!80!black,above right] at (.3,2.15) {$\gamma_2$};
 \node[below right] at (1.3,0) {$t=0$};
 \node[gray!80!black] at (-1.95,2.5) {światło};
\end{tikzpicture}
\caption{Dwie geodezyjne czasowe przechodzące przez
hiperpowierzchnię $t=0$ w czasoprzestrzeni Minkowskiego.
Nie występuje tu brzeg rozmaitości ani osobliwość
metryki; linie przerywane przedstawiają kierunki świetlne.}
\label{fig:poincare-geodezyjne-zero}
\end{figure}
\end{przyklad}

% ===== CODEX: PRZYKLAD RACHUNKOWY; ROZDZIAL KRZYWIZNY =====
% Wstaw po definicjach tensora Riemanna, Ricciego i skalara,
% po eq:riemann-wspolrzedne, przed prz:flrw-krzywizna.

\begin{przyklad}[krzywizna Poincarégo i regularność przy $t=0$]
\label{prz:poincare-krzywizna}
Dla metryki~\eqref{eq:poincare-minkowski} obliczyliśmy
w~\eqref{eq:poincare-christoffel}, że
$\Gamma^\lambda_{\mu\nu}=0$ w całej dziedzinie.
W szczególności wszystkie pochodne tych symboli
również znikają. Ze wzoru na tensor Riemanna,
zgodnie z przyjętą konwencją znaku, otrzymujemy
\begin{equation}\label{eq:poincare-riemann-zero}
 R^\lambda{}_{\rho\mu\nu}
 =\partial_\mu\Gamma^\lambda_{\nu\rho}
 -\partial_\nu\Gamma^\lambda_{\mu\rho}
 +\Gamma^\lambda_{\mu\sigma}\Gamma^\sigma_{\nu\rho}
 -\Gamma^\lambda_{\nu\sigma}\Gamma^\sigma_{\mu\rho}
 =0.
\end{equation}
Przez kontrakcję indeksów dostajemy
$\mathrm{Ric}_{\rho\nu}=R^\mu{}_{\rho\mu\nu}=0$,
a następnie $\mathrm{Scal}=g^{\rho\nu}\mathrm{Ric}_{\rho\nu}=0$.
Także tensor Einsteina
$E=\mathrm{Ric}-\frac12\mathrm{Scal}\,g_M$ oraz
niezmiennik Kretschmanna
$\mathcal K=R_{\alpha\beta\gamma\delta}
 R^{\alpha\beta\gamma\delta}$ znikają identycznie.
Są to równości tensorowe, niezależne od współrzędnych.

Co dzieje się dokładnie przy $t=0$?
Macierz metryki nadal wynosi
$\operatorname{diag}(-c^2,1,1,1)$, a jej wyznacznik
$\det(g_{\mu\nu})=-c^2\ne0$.
Hiperpowierzchnia $\Sigma_0=\{t=0\}$ jest gładka;
metryka indukowana na $\Sigma_0$ ma postać
$g_M|_{T\Sigma_0}=\dd x^2+\dd y^2+\dd z^2$.
Jej jednostkowy, przyszłościowo skierowany normalny
$U=c^{-1}\partial_t$ spełnia $g_M(U,U)=-1$.
Stosując definicję~\ref{def:izomorfizmy-muzyczne},
mamy $U^\flat=-c\,\dd t$.
Ponieważ $\nabla_XU=0$ dla każdego pola $X$
stycznego do $\Sigma_0$, druga forma fundamentalna
$\mathrm{II}(X,Y)=g_M(\nabla_XU,Y)$ także znika.
W połączeniu z geodezyjną zupełnością z
Przykładu~\ref{prz:poincare-geodezyjne} otrzymujemy,
że $t=0$ jest regularnym przekrojem płaskiej
czasoprzestrzeni, a nie osobliwością.

Podstawienie $\zeta=\mathrm{i}ct$ nie zmienia tych
wniosków: na rzeczywistym obrazie
$F(M)=\R^3\times\mathrm{i}\R$ mamy
$F^*h=g_M$ i identycznie zerowe niezmienniki.
Jest to zasadniczo inny przypadek niż pyłowa metryka
FLRW z~\eqref{eq:flrw-dust-curvature}, dla której
$\mathcal K=80/(27t^4)$ rozbiega się przy $t\to0^+$.
Różnica wynika z geometrii, nie z użycia liczby
urojonej w symbolu współrzędnej czasowej.
\end{przyklad}

% ===== CODEX: PO ISTNIEJACYCH PRZYKLADACH SCHWARZSCHILDA I KERRA =====
% ===== DO SEKCJI O METRYKACH LORENTZOWSKICH =====

\begin{przyklad}[metryka FLRW i geometria rozszerzającego się Wszechświata]
\label{prz:flrw-metryka}
Rozważmy czasoprzestrzeń $M=I\times\R^3$, gdzie $I\subset(0,\infty)$
jest przedziałem, z metryką lorentzowską
\begin{equation}\label{eq:flrw-metryka}
 g=-\dd t^2+a(t)^2(\dd x^2+\dd y^2+\dd z^2),
 \qquad a\in C^\infty(I),\quad a(t)>0.
\end{equation}
Jest to przestrzennie płaski model Friedmanna--Lemaître'a--Robertsona--Walkera
(FLRW). Funkcja $a$ nazywa się \emph{czynnikiem skali}, a
$H=\dot a/a$ parametrem Hubble'a; kropka oznacza pochodną po $t$.
Współrzędne $(x,y,z)$ są \emph{współporuszające się}.
Na przekroju $\{t\}\times\R^3$ metryka indukowana wynosi
$h_t=a(t)^2(\dd x^2+\dd y^2+\dd z^2)$, a więc jego
krzywizna wewnętrzna jest zerowa. Odległość przestrzenna dwóch
punktów o stałych współrzędnych $\mathbf x,\mathbf y$ wynosi
$D(t)=a(t)|\mathbf x-\mathbf y|$; stąd $\dot D=HD$.
Płaskość każdego przekroju przestrzennego nie oznacza jednak
płaskości czterowymiarowej czasoprzestrzeni.

Dla $x^0=t$ i indeksów przestrzennych $i,j\in\{1,2,3\}$
mamy $g_{00}=-1$, $g_{0i}=0$, $g_{ij}=a^2\delta_{ij}$
oraz $g^{00}=-1$, $g^{ij}=a^{-2}\delta^{ij}$.
W szczególności $\det(g_{\mu\nu})=-a^6$; warunek $a>0$
zapewnia niezdegenerowanie metryki na $M$.
\end{przyklad}

% ===== DO SEKCJI O KONEKSJI I SYMBOLACH CHRISTOFFELA =====

\begin{przyklad}[symbole Christoffela metryki FLRW]
\label{prz:flrw-christoffel}
Dla metryki~\eqref{eq:flrw-metryka} jedyne niezerowe
pochodne jej współczynników to
$\partial_0 g_{ij}=2a\dot a\,\delta_{ij}$.
Ze wzoru na symbole Christoffela dostajemy
\begin{equation}\label{eq:flrw-christoffel}
 \Gamma^0_{ij}=a\dot a\,\delta_{ij},\qquad
 \Gamma^i_{0j}=\Gamma^i_{j0}=H\delta^i_j,
\end{equation}
przy czym wszystkie niewypisane symbole znikają.
Istotnie,
$\Gamma^0_{ij}=\tfrac12g^{00}(-\partial_0g_{ij})
=a\dot a\,\delta_{ij}$ oraz
$\Gamma^i_{0j}=\tfrac12g^{ik}\partial_0g_{kj}
=(\dot a/a)\delta^i_j$.
Równoważnie, dla $\partial_0=\partial_t$,
\[
 \nabla_{\partial_0}\partial_0=0,\qquad
 \nabla_{\partial_0}\partial_i
 =\nabla_{\partial_i}\partial_0=H\partial_i,\qquad
 \nabla_{\partial_i}\partial_j=a\dot a\,\delta_{ij}\partial_0.
\]
\end{przyklad}

% ===== DO SEKCJI O GEODEZYJNYCH =====

\begin{przyklad}[geodezyjne czasowe i świetlne w modelu FLRW]
\label{prz:flrw-geodezyjne}
Niech $\gamma(s)=(t(s),\mathbf x(s))$ będzie afinicznie
sparametryzowaną geodezyjną metryki~\eqref{eq:flrw-metryka};
apostrof oznacza pochodną po $s$.
Podstawiając~\eqref{eq:flrw-christoffel} do równania geodezyjnej
otrzymujemy układ równań różniczkowych zwyczajnych
\begin{equation}\label{eq:flrw-ode-geodezyjnych}
 t''+a\dot a\sum_{i=1}^3(x^{i\prime})^2=0,\qquad
 x^{i\prime\prime}+2H(t)t'x^{i\prime}=0
 \quad(i=1,2,3).
\end{equation}
Drugie równanie ma całkę pierwszą
\begin{equation}\label{eq:flrw-momenty}
 \frac{\dd}{\dd s}\bigl(a(t)^2x^{i\prime}\bigr)=0,
 \qquad a(t)^2x^{i\prime}=P^i,
\end{equation}
gdzie $P^i$ są stałymi. Zatem problem geodezyjny
redukuje się do jednego równania dla $t$ i kwadratur dla $x^i$.
Dla geodezyjnej przyszłościowo skierowanej, spełniającej
$g(\gamma',\gamma')=-\varepsilon$,
gdzie $\varepsilon=1$ w przypadku czasowym (parametr $s$ jest
czasem własnym) i $\varepsilon=0$ w przypadku zerowym,
\begin{equation}\label{eq:flrw-pierwsza-calka}
 t'=\sqrt{\varepsilon+\frac{|P|^2}{a(t)^2}},\qquad
 \frac{\dd x^i}{\dd t}
 =\frac{P^i}{a(t)^2\sqrt{\varepsilon+|P|^2/a(t)^2}},
 \qquad |P|^2=\sum_i(P^i)^2.
\end{equation}
Dla $P=0$ i $\varepsilon=1$ otrzymujemy
$\gamma(t)=(t,\mathbf x_0)$: obserwatorzy o stałych
współrzędnych przestrzennych poruszają się po geodezyjnych
czasowych, a $t$ jest ich czasem własnym.

Dla geodezyjnej świetlnej $P\ne0$ wzór upraszcza się do
$\dd\mathbf x/\dd t=\mathbf n/a(t)$,
gdzie $\mathbf n=P/|P|$ jest jednostkowym kierunkiem.
Przy $a(t)=(t/t_*)^{2/3}$ całkowanie daje
\begin{equation}\label{eq:flrw-swiatlo}
 \mathbf x(t)=\mathbf x(t_e)+
 3t_*^{2/3}(t^{1/3}-t_e^{1/3})\mathbf n.
\end{equation}
Równanie to opisuje tor promienia w przestrzennych
współrzędnych współporuszających się, nie jego parametryzację afiniczną.
\end{przyklad}

\begin{przyklad}[odchylenie geodezyjnych i rozszerzanie się przestrzeni]
\label{prz:flrw-jacobi}
Dla obserwatorów współporuszających się pole prędkości wynosi
$U=\partial_t$. Dwie bliskie geodezyjne rozdziela pole Jacobiego
$J=\sum_i v^i\partial_i$ o stałych współczynnikach $v^i$.
Z~\eqref{eq:flrw-christoffel} mamy
$\nabla_UJ=HJ$ i
$\nabla_U\nabla_UJ=(\dot H+H^2)J=(\ddot a/a)J$.
Jednocześnie bezpośredni rachunek krzywizny daje
$R(J,U)U=-(\ddot a/a)J$, zatem
$\nabla_U\nabla_UJ+R(J,U)U=0$, zgodnie z równaniem Jacobiego.
Długość pola wynosi $|J|_g=a(t)|v|$ i spełnia
\[
 \frac{\dd^2}{\dd t^2}|J|_g
 =\frac{\ddot a}{a}|J|_g.
\]
Dla $\ddot a<0$ odległość nadal może rosnąć, ale tempo jej
wzrostu maleje. Równanie to wyraża wpływ krzywizny
czasoprzestrzeni na względny ruch swobodnych obserwatorów.
\end{przyklad}

% ===== DO SEKCJI O TENSORZE RIEMANNA, RICCIEGO I SKALARZE =====

\begin{przyklad}[tensor Riemanna i jego ślady dla metryki FLRW]
\label{prz:flrw-krzywizna}
Stosujemy konwencję znaku z~\eqref{eq:riemann-def}.
Ze wzorów~\eqref{eq:flrw-christoffel}
i~\eqref{eq:riemann-wspolrzedne} wynikają działania
operatora krzywizny:
\begin{align*}
 R(\partial_0,\partial_i)\partial_0
   &=\frac{\ddot a}{a}\partial_i, &
 R(\partial_0,\partial_i)\partial_j
   &=a\ddot a\,\delta_{ij}\partial_0,\\
 R(\partial_i,\partial_j)\partial_k
   &=\dot a^2(\delta_{jk}\partial_i-\delta_{ik}\partial_j).
\end{align*}
Na przykład,
$R^0{}_{i0j}=\partial_0\Gamma^0_{ji}
-\Gamma^0_{jk}\Gamma^k_{0i}
=\partial_t(a\dot a)\delta_{ij}-\dot a^2\delta_{ij}
=a\ddot a\,\delta_{ij}$.
Biorąc ślad $\mathrm{Ric}_{\mu\nu}=R^\alpha{}_{\mu\alpha\nu}$,
dostajemy
\begin{equation}\label{eq:flrw-ricci}
 \mathrm{Ric}_{00}=-3\frac{\ddot a}{a},\qquad
 \mathrm{Ric}_{0i}=0,\qquad
 \mathrm{Ric}_{ij}=(a\ddot a+2\dot a^2)\delta_{ij}.
\end{equation}
Ponieważ $g^{00}=-1$ i $g^{ij}=a^{-2}\delta^{ij}$,
krzywizna skalarna oraz tensor Einsteina
$E=\mathrm{Ric}-\tfrac12\mathrm{Scal}\,g$ przyjmują postać
\begin{align}\label{eq:flrw-scalar-einstein}
 \mathrm{Scal}
  &=6\left(\frac{\ddot a}{a}+H^2\right),\notag\\
 E_{00}&=3H^2,\qquad E_{0i}=0,\qquad
 E_{ij}=-a^2\left(2\frac{\ddot a}{a}+H^2\right)\delta_{ij}.
\end{align}

Aby odróżnić osobliwość rzeczywistą od współrzędnościowej,
wygodnie obliczyć także skalar Kretschmanna
$\mathcal K=\mathcal R_{\alpha\beta\gamma\delta}
\mathcal R^{\alpha\beta\gamma\delta}$.
W ramie ortonormalnej $e_0=\partial_t$, $e_i=a^{-1}\partial_i$
niezależne współczynniki krzywizny wyznaczają liczby
$\ddot a/a$ (kierunki czasoprzestrzenne) i $H^2$
(kierunki czysto przestrzenne), skąd
\begin{equation}\label{eq:flrw-kretschmann}
 \mathcal K=12\left[\left(\frac{\ddot a}{a}\right)^2+H^4\right].
\end{equation}
W szczególności przekroje $t=\mathrm{const}$ są płaskie,
choć zwykle $\mathcal K\ne0$.
\end{przyklad}

% ===== DO SEKCJI O ZASTOSOWANIACH KRZYWIZNY / ROWNANIACH EINSTEINA =====

\begin{przyklad}[równania Einsteina i wyprowadzenie równań Friedmanna]
\label{prz:flrw-friedmann}
Przyjmujemy $c=1$, oznaczamy stałą Newtona przez $G_{\!N}$
i rozważamy równania Einsteina z członem kosmologicznym
\begin{equation}\label{eq:flrw-einstein}
 \mathrm{Ric}_{\mu\nu}-\frac12\mathrm{Scal}\,g_{\mu\nu}
 +\Lambda g_{\mu\nu}=8\pi G_{\!N}T_{\mu\nu}.
\end{equation}
Zakładamy, że materia jest płynem doskonałym
o gęstości energii $\rho(t)$, ciśnieniu $p(t)$
i czteroprędkości $U=\partial_t$:
$T=(\rho+p)U^\flat\otimes U^\flat+pg$.
W rozważanych współrzędnych $T_{00}=\rho$,
$T_{0i}=0$, $T_{ij}=pa^2\delta_{ij}$.
Składowa $(0,0)$ równania~\eqref{eq:flrw-einstein}
oraz składowe przestrzenne, po użyciu
\eqref{eq:flrw-scalar-einstein}, dają odpowiednio
\begin{equation}\label{eq:flrw-friedmann}
 H^2=\frac{8\pi G_{\!N}}3\rho+\frac{\Lambda}{3},
 \qquad
 \frac{\ddot a}{a}=-\frac{4\pi G_{\!N}}3(\rho+3p)
 +\frac{\Lambda}{3}.
\end{equation}
Są to \emph{równania Friedmanna} dla przestrzennie płaskiego
modelu FLRW. Równoważnie drugie z nich, razem z pierwszym,
można zapisać jako
$2\ddot a/a+H^2=\Lambda-8\pi G_{\!N}p$.
Tożsamość Bianchiego i~\eqref{eq:flrw-einstein} implikują
$\nabla_\mu T^{\mu\nu}=0$. Dla $\nu=0$ jest to równanie
ciągłości
\begin{equation}\label{eq:flrw-ciaglosc}
 \dot\rho+3H(\rho+p)=0.
\end{equation}
Do wyznaczenia $a(t)$ potrzeba jeszcze \emph{równania stanu}
łączącego $p$ i $\rho$: same równania Einsteina
nie określają składu materii.
\end{przyklad}

\begin{przyklad}[rozwiązanie równań Friedmanna: pył i promieniowanie]
\label{prz:flrw-rozwiazanie-ode}
Niech $\Lambda=0$, $\rho>0$, a równanie stanu ma postać
$p=w\rho$ ze stałą $w>-1$. Równanie~\eqref{eq:flrw-ciaglosc}
daje po rozdzieleniu zmiennych
\[
 \frac{\dd\rho}{\rho}=-3(1+w)\frac{\dd a}{a},
 \qquad \rho=\rho_0\left(\frac{a_0}{a}\right)^{3(1+w)}.
\]
Podstawiając do pierwszego równania~\eqref{eq:flrw-friedmann}
i wybierając gałąź rozszerzającą się ($\dot a>0$), otrzymujemy
\[
 \dot a=C a^{-(1+3w)/2},\qquad
 C=\sqrt{\frac{8\pi G_{\!N}}{3}\rho_0a_0^{3(1+w)}}.
\]
To równanie różniczkowe zwyczajne jest rozdzielne.
Całkowanie daje
$\frac{2}{3(1+w)}a^{3(1+w)/2}=C(t-t_B)$,
a zatem
\begin{equation}\label{eq:flrw-power-law}
 a(t)=A(t-t_B)^{\frac{2}{3(1+w)}},\qquad
 \rho(t)=\frac{1}{6\pi G_{\!N}(1+w)^2(t-t_B)^2},
 \quad t>t_B,
\end{equation}
gdzie $A>0$ i $t_B$ są stałymi.
Dla \emph{pyłu} ($p=0$, $w=0$), po wyborze początku czasu
$t_B=0$ i normalizacji $a(t_*)=1$, mamy
\begin{equation}\label{eq:flrw-dust}
 a(t)=\left(\frac{t}{t_*}\right)^{2/3},\qquad
 H(t)=\frac{2}{3t},\qquad
 \rho(t)=\frac{1}{6\pi G_{\!N}t^2}.
\end{equation}
Dla \emph{promieniowania} ($p=\rho/3$, $w=1/3$)
otrzymujemy $a(t)=(t/t_*)^{1/2}$ oraz
$\rho(t)=3/(32\pi G_{\!N}t^2)$.
Natomiast przy $w=-1$ gęstość jest stała, a równanie
$\dot a=Ha$ daje $a(t)=a_0e^{H(t-t_0)}$,
$H^2=8\pi G_{\!N}\rho/3$; w tym przypadku
nie występuje osobliwość krzywizny przy skończonym $t$.

\begin{figure}[htbp]
\centering
\begin{tikzpicture}[>=Stealth,x=1.05cm,y=.95cm]
 \draw[->] (0,0)--(4.8,0) node[right] {$t$};
 \draw[->] (0,0)--(0,3.1) node[above] {$a(t)$};
 \draw[very thick,manifoldblue,domain=0:4,samples=80]
  plot (\x,{2.5*pow(\x/4,0.6666667)});
 \draw[thick,manifoldred,dashed,domain=0:4,samples=80]
  plot (\x,{2.5*sqrt(\x/4)});
 \draw[densely dotted] (4,0)--(4,2.5);
 \fill (4,2.5) circle (1.4pt);
 \node[below] at (4,0) {$t_*$};
 \node[left] at (0,2.5) {$1$};
 \node[manifoldblue,anchor=west] at (1.95,1.25) {$w=0$};
 \node[manifoldred,anchor=west] at (1.05,2.4) {$w=\frac13$};
 \fill (0,0) circle (1.3pt);
\end{tikzpicture}
\caption{Czynnik skali dla pyłu (linia ciągła) i promieniowania
(linia przerywana), przy $a(t_*)=1$. Punkt $t=0$ jest granicą
wykresów, a nie punktem rozważanej czasoprzestrzeni.}
\label{fig:flrw-skala}
\end{figure}
\end{przyklad}

\begin{przyklad}[czas konforemny i geodezyjne świetlne]
\label{prz:flrw-konforemny}
Dla modelu pyłowego~\eqref{eq:flrw-dust} wprowadźmy
\emph{czas konforemny} $\eta$ przez $\dd\eta=\dd t/a(t)$.
Wybierając $\eta\to0$ dla $t\to0^+$, otrzymujemy
$\eta=3t_*^{2/3}t^{1/3}$ oraz
\begin{equation}\label{eq:flrw-konforemna}
 g=a(\eta)^2\bigl(-\dd\eta^2+\dd x^2+\dd y^2+\dd z^2\bigr).
\end{equation}
Ponieważ metryka jest konforemna do płaskiej metryki
Minkowskiego, jej tensor Weyla znika identycznie.
Promienie świetlne, rozumiane jako \emph{niesparametryzowane}
geodezyjne zerowe, są zatem prostymi
$\mathbf x(\eta)=\mathbf x_e+\mathbf n(\eta-\eta_e)$,
$|\mathbf n|=1$. To nie znaczy, że $\eta$ jest ich
parametrem afinicznym: dla niezerowego pędu $P$
z~\eqref{eq:flrw-momenty} mamy
$\dd\lambda/\dd t=a(t)/|P|$.

\begin{figure}[htbp]
\centering
\begin{tikzpicture}[>=Stealth,x=1.05cm,y=1.05cm]
 \draw[->] (-2.6,0)--(2.75,0) node[right] {$x$};
 \draw[->] (0,0)--(0,3.6) node[above] {$\eta$};
 \draw[densely dashed,manifoldred] (-2.5,0.08)--(2.5,0.08);
 \node[manifoldred,below left] at (-1.6,0.08) {$t=0$};
 \foreach \z in {-2,-1,1,2} {
  \draw[gray!65] (\z,0.16)--(\z,3.3);
 }
 \node[gray!75!black] at (1.8,3.48) {obserwatorzy współporuszający się};
 \fill[manifoldblue] (0,.75) circle (2pt);
 \draw[very thick,manifoldblue,->] (0,.75)--(-2.1,2.85);
 \draw[very thick,manifoldblue,->] (0,.75)--(2.1,2.85);
 \node[manifoldblue,right] at (1.7,2.5) {światło};
\end{tikzpicture}
\caption{Przekrój diagramu konforemnego FLRW.
Geodezyjne współporuszające się są pionowe, a promienie świetlne
mają nachylenie $45^\circ$. Linia $t=0$ jest brzegiem
wybranej czasoprzestrzeni, nie jej przekrojem.}
\label{fig:flrw-konforemny}
\end{figure}
\end{przyklad}

\begin{przyklad}[osobliwość Wielkiego Wybuchu: krzywizna i niezupełność]
\label{prz:flrw-osobliwosc}
W modelu pyłowym~\eqref{eq:flrw-dust} mamy
$\dot a/a=2/(3t)$ i $\ddot a/a=-2/(9t^2)$.
Wstawienie do~\eqref{eq:flrw-ricci}--\eqref{eq:flrw-kretschmann}
daje
\begin{align}\label{eq:flrw-dust-curvature}
 \mathrm{Ric}_{00}&=\frac{2}{3t^2},&
 \mathrm{Ric}_{ij}&=\frac{2a(t)^2}{3t^2}\delta_{ij},\notag\\
 \mathrm{Scal}&=\frac{4}{3t^2},&
 \mathcal K&=\frac{80}{27t^4}.
\end{align}
Zatem $\mathcal K\to+\infty$ przy $t\to0^+$.
Nie jest to osobliwość wynikająca wyłącznie ze złego wyboru
współrzędnych: $\mathcal K$ jest niezmiennikiem skalarnym,
więc żadna zmiana współrzędnych nie może usunąć jego
rozbieżności. W szczególności nie istnieje przedłużenie
tego rozwiązania przez $t=0$ jako regularnej metryki
lorentzowskiej klasy $C^2$, przez którą przechodzą te krzywe.

Ten sam problem widać na geodezyjnych.
Obserwator współporuszający się ma $\dd\tau=\dd t$,
więc od $t=t_0$ do granicy $t=0$ upływa skończony
czas własny $\Delta\tau=t_0$.
Dla geodezyjnej świetlnej z $P\ne0$ wzory
\eqref{eq:flrw-momenty}--\eqref{eq:flrw-pierwsza-calka} dają
\[
 \Delta\lambda=\frac1{|P|}\int_0^{t_0}a(t)\,\dd t
 =\frac{3t_0^{5/3}}{5|P|t_*^{2/3}}<\infty.
\]
Geodezyjne są więc \emph{niezupełne ku przeszłości}:
kończą się po skończonym parametrze afinicznym, choć
$M=(0,\infty)\times\R^3$ nie zawiera punktów o $t=0$.

Warto odróżnić brak wartości wzorów w $t=0$ od dokładnego
wniosku geometrycznego. Dla rozważanego \emph{klasycznego
rozwiązania} równań Einsteina z pyłem nie ma regularnego
przedłużenia przez tę granicę; nie wynika stąd, że
\emph{każde} rozwiązanie równań Einsteina ma osobliwość
w $t=0$ albo że teoria nie opisuje czasów $t>0$.
\end{przyklad}

% ===== CODEX: ZARAZ PO prz:flrw-osobliwosc, PRZED prz:flrw-milne =====

\begin{przyklad}[czas zespolony a osobliwość w modelu FLRW]
\label{prz:flrw-czas-zespolony}
Sprawdźmy, czy przedstawienie Poincarégo z
Przykładu~\ref{prz:poincare-czas-urojony} pomaga
zrozumieć granicę $t\to0^+$ modelu pyłowego.
Dla metryki~\eqref{eq:flrw-metryka} z~\eqref{eq:flrw-dust},
\[
 g=-\dd t^2+\left(\frac{t}{t_*}\right)^{4/3}
    (\dd x^2+\dd y^2+\dd z^2),\qquad t>0.
\]
Po formalnym podstawieniu $\zeta=\mathrm{i}t$,
czyli $t=-\mathrm{i}\zeta$, otrzymujemy
\begin{equation}\label{eq:flrw-complex-time}
 g=\dd\zeta^2+
 \left(\frac{-\mathrm{i}\zeta}{t_*}\right)^{4/3}
 (\dd x^2+\dd y^2+\dd z^2).
\end{equation}
Na obrazie rzeczywistej półprostej $t>0$ mamy
$\zeta\in\mathrm{i}(0,\infty)$; zapis ten nadal
reprezentuje tę samą metrykę lorentzowską.
Aby wyjść poza tę półprostą w płaszczyźnie
zespolonej, trzeba wybrać gałąź funkcji
$z^{4/3}$; punkt $z=0$ jest punktem rozgałęzienia.

Nie jest to automatycznie \emph{obrót Wicka}.
W tym drugim zabiegu próbuje się analitycznie
kontynuować rozwiązanie do $t=-\mathrm{i}\tau$,
$\tau\in\R$, a następnie otrzymać rzeczywistą
metrykę riemannowską. W naszym przypadku dla
$\tau>0$ na gałęzi głównej dostajemy jedynie
zespoloną formę
\[
 g_{\mathrm E}^{\mathrm{formal}}
 =\dd\tau^2+e^{-2\pi\mathrm{i}/3}
  \left(\frac{\tau}{t_*}\right)^{4/3}
  (\dd x^2+\dd y^2+\dd z^2),
\]
która nie jest rzeczywistą metryką riemannowską.

Co ważniejsze, kontynuowanie samego niezmiennika
Kretschmanna z~\eqref{eq:flrw-dust-curvature}
daje
\[
 \mathcal K(-\mathrm{i}\zeta)
  =\frac{80}{27\zeta^4},\qquad
 \mathcal K(-\mathrm{i}\tau)
  =\frac{80}{27\tau^4}.
\]
W pierwszym wzorze $\zeta=\mathrm{i}t$ i stosujemy
$\mathcal K(t)=80/(27t^4)$, a w drugim
kontynuujemy bezpośrednio $t\mapsto-\mathrm{i}\tau$.
W obu zapisach biegun czwartego rzędu pozostaje
w początku. Kontynuacja formalnego wzoru na
niezmiennik nie zapewnia wprawdzie istnienia
rzeczywistej metryki, ale wyklucza usunięcie
tej rozbieżności przez samo przemianowanie czasu.

\begin{figure}[htbp]
\centering
\begin{tikzpicture}[>=Stealth,x=1.1cm,y=1.1cm]
 \draw[->] (-2.45,0)--(3.15,0) node[right] {$\operatorname{Re}t$};
 \draw[->] (0,-2.4)--(0,2.35) node[above] {$\operatorname{Im}t$};
 \draw[manifoldblue,very thick,->] (.16,0)--(2.8,0);
 \node[manifoldblue,above] at (2.13,.06) {$t>0$};
 \draw[manifoldred,very thick,->] (0,-.13)--(0,-2.12);
 \node[manifoldred,right] at (.07,-1.86) {$t=-\mathrm{i}\tau$};
 \draw[manifoldgold!85!black,very thick,->,
       domain=0:80,samples=40,variable=\a]
   plot ({1.35*cos(-\a)},{1.35*sin(-\a)});
 \node[manifoldgold!70!black,right] at (1.38,-1.05)
   {kontynuacja formalna};
 \fill[white] (0,0) circle (3pt);
 \draw[manifoldred,very thick] (0,0) circle (3pt);
 \node[manifoldred,above left] at (-.08,.11)
   {$t=0$};
\end{tikzpicture}
\caption{Schemat płaszczyzny czasu zespolonego.
Dodatnia półprosta rzeczywista odpowiada fizycznemu
czasowi modelu FLRW, a półprosta ujemnie urojona
ilustruje próbę obrotu Wicka. Punkt $t=0$
pozostaje osobliwy; łuk przedstawia kontynuację
formalną, nie izometrię rzeczywistych czasoprzestrzeni.}
\label{fig:flrw-wick}
\end{figure}

Nie twierdzimy, że każde rozwiązanie równań
Einsteina jest osobliwe. Wniosek dotyczy tej
konkretnej metryki pyłowej oraz jej rozbieżnych
niezmienników krzywizny. W innych geometriach
można uzyskać regularne kontynuacje euklidesowe,
ale wymaga to osobnej konstrukcji.
\end{przyklad}

\begin{przyklad}[dlaczego zanik czynnika skali nie wystarcza]
\label{prz:flrw-milne}
Dla porównania rozważmy na $t>0$ metrykę
\[
 g=-\dd t^2+t^2\bigl(\dd\chi^2+
 \sinh^2\chi\,(\dd\theta^2+\sin^2\theta\,\dd\phi^2)\bigr).
\]
Ma ona czynnik skali $a(t)=t$ i ujemną krzywiznę
przekrojów przestrzennych (model Milne'a).
Podstawienie $T=t\cosh\chi$, $R=t\sinh\chi$ daje
\[
 g=-\dd T^2+\dd R^2+
 R^2(\dd\theta^2+\sin^2\theta\,\dd\phi^2),
\]
czyli metrykę Minkowskiego we wnętrzu przyszłego stożka
świetlnego. Tensor Riemanna znika identycznie;
pozorny problem przy $a(0)=0$ wynika tutaj z niepełnego
pokrycia współrzędnymi, a nie z rozbieżności krzywizny.
Wzór~\eqref{eq:flrw-kretschmann} dotyczył wyłącznie
przestrzennie płaskiego modelu FLRW i nie wolno go
stosować do tej metryki.

Jeszcze inny kontrast stanowi płaski przestrzennie model
z promieniowaniem: $a(t)\sim t^{1/2}$, dla którego
$\mathrm{Scal}=0$, ale $\mathcal K=3/(2t^4)\to\infty$.
Sam brak rozbieżności krzywizny skalarnej nie świadczy
zatem o regularności czasoprzestrzeni.
\end{przyklad}
